% V7 evidence-bounded Method draft. Not submission ready.
\section{Reliability-Guided GPS Covariance Adaptation}
\label{sec:method_v7}

\subsection{System Scope and Batch Factor-Graph Formulation}
The implemented estimator optimizes vehicle poses
$\mathbf{T}_{0:N-1}\in\mathrm{SE}(3)$ using synthetically degraded
GPS positions, gyroscope measurements, LiDAR odometry, and stereo
visual odometry. The formulation is an \emph{offline batch} pose graph,
not an online incremental smoother or a full velocity--bias
state estimator. A fixed identity prior anchors the first pose.
A GPS factor constrains translation at each time step, while
relative gyroscope, LiDAR, and stereo constraints couple adjacent
poses. The GPS factor uses effectively uninformative rotational
noise, and the three non-GPS noise settings are unchanged between
compared methods.

For illustration, the optimized objective is
\begin{equation}
\begin{aligned}
\min_{\mathbf{T}_{0:N-1}}\quad&
\|\mathbf{r}_{0}\|_{\mathbf{\Omega}_0}^{2}+
\sum_{t=0}^{N-1}
\|\mathbf{p}(\mathbf{T}_t)-\mathbf{z}^{g}_t
\|_{\mathbf{\Sigma}^{g}_t}^{2}\\
&+\sum_{t=0}^{N-2}\bigl(
\|\mathbf{r}^{i}_{t,t+1}\|_{\mathbf{\Omega}_i}^{2}+
\|\mathbf{r}^{l}_{t,t+1}\|_{\mathbf{\Omega}_l}^{2}+
\|\mathbf{r}^{c}_{t,t+1}\|_{\mathbf{\Omega}_c}^{2}\bigr),
\end{aligned}
\label{eq:v7_objective}
\end{equation}
where $\|\mathbf{r}\|_{\mathbf{\Sigma}}^2$ denotes Mahalanobis
weighting and $\mathbf{p}(\mathbf{T})$ extracts position.
The implementation uses GTSAM \texttt{PriorFactorPose3} and
\texttt{BetweenFactorPose3}. The inertial link is a
gyroscope-derived relative-rotation pseudo-factor; it is not
a full preintegrated IMU factor. The graph is optimized by
Levenberg--Marquardt.

\subsection{Causal Features and Reliability Networks}
The reliability model observes only present and historical
\emph{degraded} GPS positions. Its 21-dimensional feature at
each frame consists of position increments, nominal
finite-difference velocity and acceleration, and sliding-window
statistics of their magnitudes. The nominal sampling period is
$0.1$ s and the feature-statistics window spans ten frames.
A 64-frame causal sequence is normalized by training-only
feature medians and interquartile ranges, then clipped to
$[-20,20]$.

Two-layer GRU encoders are compared with a two-block Mamba encoder.
A GRU with width 96 contains 114,630 trainable parameters;
a width-106 GRU capacity control contains 139,290; Mamba with
width 96 contains 139,590. All use a shared input-projection
design, comparable Beta output heads, and the same supervision
objective. The Mamba case uses actual Mamba sequence blocks
rather than a recurrent substitute.

Each network predicts Beta-distribution parameters for
$h\in\{0,1,2\}$, representing current and two future
synthetic reliability targets:
\begin{equation}
\alpha_{t,h}=1+\mathrm{softplus}(a_{t,h})+\epsilon,\qquad
\beta_{t,h}=1+\mathrm{softplus}(b_{t,h})+\epsilon,
\end{equation}
\begin{equation}
\hat r_{t,h}=
\frac{\alpha_{t,h}}{\alpha_{t,h}+\beta_{t,h}},\qquad
U_{t,h}=\frac{2}{\alpha_{t,h}+\beta_{t,h}}.
\label{eq:v7_reliability}
\end{equation}
Model training minimizes a weighted Beta negative log likelihood
for the three heads with relative weights $1.0:0.7:0.5$, plus
a weak evidence penalty. Training and model selection use the
already explored KITTI 0027 development data; the nominally
held-out drives do not provide network training signals.
The quantity $U$ is a Beta concentration proxy, not a validated
epistemic uncertainty estimate or physically calibrated noise.

\subsection{Training-Only Axis Noise Mapping and Cap4}
A five-bin axis-noise mapping is fitted from 0027
\emph{training} degraded GPS errors and synthetic oracle
reliability targets. For reliability bin $b$ and
axis $a\in\{x,y,z\}$, the empirical RMS error is
\begin{equation}
s_{b,a}=\sqrt{\frac{1}{|\mathcal I_b|}
\sum_{t\in\mathcal I_b}
(z^{\mathrm{deg}}_{t,a}-z^{\mathrm{ref}}_{t,a})^2}.
\label{eq:v7_axisfit}
\end{equation}
The original implementation applies monotone enforcement
and a positive floor to the bin values.
Linear interpolation using the predicted current reliability
$\hat r_{t,0}$ yields $\tilde s_{t,a}$. To limit
overconfident weighting, the dynamic axis sigma is
\begin{equation}
s_{t,a}=\max\left(
\tilde s_{t,a},
s_a^{\mathrm{fixed}}/\sqrt{4}\right),\quad
\mathbf{\Sigma}^{g}_t=
\mathrm{diag}(s_{t,x}^2,s_{t,y}^2,s_{t,z}^2).
\label{eq:v7_cap4}
\end{equation}
Thus GPS precision cannot exceed four times the
training-only fixed-axis precision. This Cap4
restriction is an engineering heuristic and is
\emph{not} a proven probabilistic calibration.
For the first 63 frames, all cases use the same
training-only fixed-axis noise because a full
64-frame context is unavailable.

Importantly, the tested graph uses only
$\hat r_{t,0}$ for weighting. The $h=1,2$
prediction heads and evidence proxy $U$ are
\emph{not} fed into the evaluated factor graph.
Consequently, these experiments do not establish
anticipatory reliability weighting, feedforward--feedback
control, factor forgetting, or iSAM2 performance.

\subsection{Optimization Consistency}
For a given drive, every comparison uses identical
generated GPS, IMU, LiDAR, and stereo factors,
date-specific calibration, and identical
LiDAR-integrated initial poses. No optimized
development trajectory is reused to initialize
the held-out drives. Every case uses at most
12 LM passes with at most 100 iterations per pass.
A run is marked numerically stable only when
relative objective reduction is at most $10^{-6}$
and the RMS position update is at most $0.01$ m.
Numerical stability does not certify a global optimum.
